% Generated draft: requires mathematical and source review.
\documentclass[11pt]{article}
\usepackage[margin=25mm]{geometry}
\usepackage{fontspec}
\setmainfont{Latin Modern Roman}
\usepackage{amsmath,amssymb,mathtools,unicode-math}
\setmathfont{Latin Modern Math}
\usepackage{enumitem,needspace,xurl,hyperref}
\setlength{\parindent}{0pt}
\setlength{\parskip}{5pt}
\setlength{\emergencystretch}{3em}
\newcommand{\N}{\mathbb{N}}
\newcommand{\Z}{\mathbb{Z}}
\newcommand{\Q}{\mathbb{Q}}
\newcommand{\R}{\mathbb{R}}
\newcommand{\C}{\mathbb{C}}
\newcommand{\F}{\mathbb{F}}
\newcommand{\A}{\mathbb{A}}
\newcommand{\PP}{\mathbb{P}}
\newcommand{\E}{\mathbb{E}}
\newcommand{\eps}{\varepsilon}
\newcommand{\dd}{\,\mathrm{d}}
\begin{document}
\section*{UnsolvedMath ID: 20000551}
\textbf{Source problem code:} catalogue entry
\subsection{Definitions and mathematical statement}
The original bullet asks whether ``Pete's example'' can be improved, saying it affects the equivalences and is almost a counterexample. The source does not define Pete, the example, target inequality, or meaning of better. Its verification says neighboring bullets discuss matching quantitative lower bounds, an order-\(\varepsilon^{-2}\) real-paving bound, and a constructive bound of ``Pete C.'' No unique original mathematical problem is recoverable.

The report explicitly adopts one interpretation using the Casazza--Fickus--Mixon--Tremain two-block DFT projection. For that particular construction, it proposes that every first-block principal compression indexed by \(S\), \(|S|=s\geq n\), has
\[
P[S]=\frac1{n+1}I+K_S,\qquad K_S\geq0,\qquad\operatorname{rank}K_S=n-1.
\]
Consequently its eigenvalue \(1/(n+1)\) has exact multiplicity \(s-n+1\). The defining DFT matrices, projection normalization, complete index range, and admissible \(n\) are not supplied in the summary, so this is not a statement about arbitrary projections.

For every bipartition, the report asserts that the two complementary-projection compression spectra have together at least two eigenvalues at least \(n/(n+1)\). In the natural paving interpretation these are the two block compressions of \(I-P\); the full index convention is not printed. It further says one block of the associated zero-diagonal reflection has norm at least
\[
\frac{n-1}{n+1}.
\]
A reflection is represented by \(2P-I\), up to overall sign, in the constant-half-diagonal projection convention. These normalization qualifications are essential to the claimed spectral consequence.

The report additionally claims a Weyl-stability corollary and an exact even/odd replication no-amplification lemma, but gives neither their hypotheses nor formulas. It separately identifies a conference-matrix reading of the ambiguous workshop remark and derives its known order-\(\varepsilon^{-2}\) paving lower bound. This alternative reading is not identified as the original example. The results are candidate refinements of specified surrogate constructions, not a recovery or resolution of the undefined ``make it better'' request.
\subsection{Short English statement}
The workshop example is unidentified. Under an explicit DFT-projection interpretation, the report proposes an exact small-eigenvalue multiplicity and consequent near-unpaved modes, while retaining a separate conference-matrix interpretation and missing construction details.
\subsection{Sources}
\begin{itemize}
\item AIM workshop ``Beyond Kadison--Singer: paving and consequences,'' December 1--5, 2014, item b10: \url{https://docs.google.com/document/d/1XRw0WyG5PshqboCp99tovCRD1GCK0jXZiNLH_nw618c/edit}. The supplied verification places the bullet under a list headed by Petter Br\"and\'en; the source document was not accessed here.
\item Local canonical registry \texttt{aim-analysis-notes.json}, \texttt{notes[82]}, tag \texttt{question}; permanent UnsolvedMath ID \(20000551\), source code \texttt{catalogue entry}. The original only refers to Pete's example. The verification gives neighboring ``Pete C.'' and quantitative-paving context, but explicitly says no unique statement is recoverable. Original remarks and literature field are empty; the clean statement is absent.
\item Supplied research attempt 1 adopts the Casazza--Fickus--Mixon--Tremain DFT construction and proposes spectral refinements. No authoritative full references, matrices, or exact replication/stability statements are supplied. The report is machine generated, not peer reviewed, and requires independent verification before citation.
\end{itemize}
\end{document}
